% ===================================================================
%  TAVOLA N. 13 — L'albergo. --- Il ristorante. --- Il caffè.
%  Traduction 2026 d'apres texto/io, controlee sur texto/fr et
%  texto/en. Consigne : voir texto/it/10-tavola-01.tex.
%
%  Deux notes marquees toutes deux « (*) », et deux appels : celui du
%  bloc 01-2 (la chambre) et celui du bloc 09-1 (le menu).
%
%  LE POURBOIRE DE L'ALINEA 9 EST UNE MANCHE DE ROBE.
%
%  Le tableau 13 allemand avait compare les mots du pourboire :
%  allemand Trinkgeld, l'argent-a-boire ; estonien
%  jootraha, l'argent-a-boire, calque de l'allemand ;
%  francais pourboire, la meme image encore ; anglais
%  tip, qui ne la partage pas. Trois langues sur quatre
%  paient a boire.
%
%  L'italien dit mancia, de l'ancien francais « manche » : LA
%  MANCHE que la dame detachait de sa robe et donnait au chevalier
%  comme faveur au tournoi. Le mot a glisse du gage courtois au
%  cadeau, puis du cadeau au pourboire du garcon.
%
%     QUATRE LANGUES REGARDENT LE MEME GESTE ET N'Y VOIENT PAS LA MEME
%     CHOSE : trois y voient un verre a payer, l'anglais un petit
%     coup, et l'italien un present de tournoi. LE MOT NE DECRIT PAS
%     LE GESTE, IL DECRIT L'ENDROIT D'OU ON LE REGARDE.
%
%  ET L'ALBERGO DU TITRE EST UN CANTONNEMENT MILITAIRE. « Albergo »
%  vient du germanique « *haribergo » — « hari », l'armee, et
%  « bergan », abriter : L'ABRI DE LA TROUPE. C'est le meme mot que
%  l'« auberge » francaise et la « Herberge » allemande, arrive en
%  Italie avec les Lombards. L'Hotel Oceano de l'alinea 1
%  porte, lui, la forme francaise de « hospitale », l'hopital : les
%  deux mots de la page pour loger un homme sont l'un une caserne,
%  l'autre un hospice.
%
%  ENFIN LE CAFFÈ DE L'ALINEA 12 EST DEUX CHOSES SOUS UN SEUL MOT.
%  L'italien appelle « caffè » la boisson et l'etablissement, sans
%  jamais les distinguer. Toutes les autres langues ont du les
%  separer : Kaffee et Café, coffee et café, кофе et кафе. L'italien
%  ne l'a pas fait parce qu'il n'en a jamais eu besoin — le premier
%  cafe d'Europe a ouvert a Venise en 1645, et le lieu s'appelait par
%  ce qu'on y buvait.
%
%     L'AUBERGE EST NOMMEE PAR CELUI QUI Y DORMAIT, LE CAFE PAR CE
%     QU'ON Y BOIT. L'un regarde l'hote, l'autre la tasse.
%
%  Le tableau 13 allemand avait trouve un hotel qui garde ses mots
%  francais par prestige ; ici il n'y a rien a garder, parce que ces
%  mots-la sont deja les siens ou ceux de ses voisins immediats.
% ===================================================================

%%K t13-noto-1 noto
\VUnotes{143.2mm}{%
\textit{\VUgras{(*) Per i particolari della camera vedi la tavola n. 6.}}%
}

%%K t13-apar-1 apar
\VUsaut{3.0mm}
\VUtitre{15.0pt}{60}{TAVOLA N. 13}
\VUsaut{1.6mm}
\VUpk{10.2pt}{40}{L'albergo. --- Il ristorante.}
\VUsaut{2.0mm}
\VUpk{10.2pt}{40}{Il caffè.}
\VUsaut{3.4mm}

%%K t13-01-1 p
1. --- Sono arrivato a Royan ieri sera e ho preso alloggio all'\textit{Hotel Oceano},
che mi aveva consigliato un amico.

%%K t13-01-2 p
Mi hanno dato una bella \VUgras{camera} \textsuperscript{(2)} al secondo
\VUgras{piano} \textsuperscript{(1)}, dove ho dormito benissimo (*).

%%K t13-02-1 p
2. --- Stamattina, uscendo dalla mia camera, dove la \VUgras{cameriera}
\textsuperscript{(3)} mi aveva portato la colazione, sono andato nella \VUgras{sala da
fumo} \textsuperscript{(5)}, che serve anche da \VUgras{sala di lettura}
\textsuperscript{(4)}, a sfogliare i \VUgras{giornali} \textsuperscript{(6)} fumando
una \VUgras{sigaretta} \textsuperscript{(8)}. Finito di leggere, sono sceso con
l'\VUgras{ascensore} \textsuperscript{(9)} e sono andato a fare un giro per la città.

%%K t13-03-1 p
3. --- Siccome avevo dimenticato di chiedere all'\VUgras{impiegato}
\textsuperscript{(12)} dell'albergo a che ora si mangia alla \VUgras{tavola rotonda}
\textsuperscript{(11)}, sono tornato presto e ne ho approfittato per fare un bagno
nella \VUgras{stanza da bagno} \textsuperscript{(10)} dell'albergo. Poi sono andato
alla \VUgras{cassa} \textsuperscript{(15)} per telefonare a un amico. Ma il
\VUgras{telefono} \textsuperscript{(13)} era occupato; vedendo che ero in difficoltà,
la \VUgras{cassiera} \textsuperscript{(14)}, che stava registrando un \VUgras{conto}
\textsuperscript{(17)} sul \VUgras{registro dei clienti} \textsuperscript{(16)}, ha
chiesto al \VUgras{portiere} \textsuperscript{(18)} di mandare il \VUgras{fattorino}
\textsuperscript{(19)} a fare la mia commissione con la sua \VUgras{bicicletta}
\textsuperscript{(20)}.

%%K t13-04-1 p
4. --- L'ho ringraziata e sono uscito a dare la mancia al \VUgras{facchino}
\textsuperscript{(24)} (l'uomo di fatica), che sul suo \VUgras{carretto}
\textsuperscript{(25)} mi aveva portato dalla stazione la mia \VUgras{cappelliera}
\textsuperscript{(26)}, la mia \VUgras{valigia} \textsuperscript{(27)} e la mia
\VUgras{borsa da viaggio} \textsuperscript{(28)}. Il \VUgras{postino}
\textsuperscript{(21)}, che arrivava proprio allora, mi ha consegnato una
\VUgras{lettera} \textsuperscript{(22)}.

%%K t13-05-1 p
5. --- Sono rimasto ancora un po' sulla soglia, vicino alla \VUgras{buca delle
lettere} \textsuperscript{(23)}, a guardare il movimento della strada. Il
\VUgras{cocchiere} \textsuperscript{(30)} dell'\VUgras{omnibus} \textsuperscript{(29)}
stava caricando sulla sua vettura il \VUgras{baule} \textsuperscript{(31)} di un
\VUgras{cliente} \textsuperscript{(32)} a cui l'\VUgras{albergatore}
\textsuperscript{(33)} diceva addio. Un giovane \VUgras{fattorino del telegrafo}
\textsuperscript{(35)} portava, senza nessuna fretta, un \VUgras{telegramma}
\textsuperscript{(34)} per uno dei clienti dell'albergo; un \VUgras{turista}
\textsuperscript{(36)}, con la sua \VUgras{macchina fotografica}
\textsuperscript{(38)} a tracolla, consultava la sua \VUgras{guida}
\textsuperscript{(37)}.

%%K t13-tit-1 sub
\VUsaut{4.0mm}
\VUpk{10.2pt}{40}{L'ora di pranzo.}
\VUsaut{3.0mm}

%%K t13-06-1 p
6. --- Davanti alla cancellata un placido \VUgras{portabagagli}
\textsuperscript{(39)} sonnecchiava fra i suoi \VUgras{ganci da carico}
\textsuperscript{(40)} e la sua \VUgras{cassetta da lustrascarpe}
\textsuperscript{(41)}, mentre una giovane \VUgras{lavandaia} \textsuperscript{(42)}
con un \VUgras{cesto} \textsuperscript{(43)} di biancheria rideva dell'aria solenne del
\VUgras{capocameriere} \textsuperscript{(44)} fermo sulla porta.

%%K t13-07-1 p
7. --- La vista del \VUgras{cuoco} \textsuperscript{(45)}, uscito dalla sua
\VUgras{cucina} \textsuperscript{(47)} nel \VUgras{seminterrato}
\textsuperscript{(48)} per mandare uno \VUgras{sguattero} \textsuperscript{(46)} a fare
una commissione, mi ha ricordato che era quasi ora di pranzo, e sono andato al
ristorante attraverso il cortile, dove un \VUgras{automobilista} \textsuperscript{(50)}
stava posteggiando la sua \VUgras{automobile} \textsuperscript{(49)} mentre un
\VUgras{meccanico} \textsuperscript{(52)} riparava, seguendo le indicazioni del
\VUgras{ciclista} \textsuperscript{(53)}, un \VUgras{triciclo a motore}
\textsuperscript{(51)} che aveva appena avuto un incidente.

%%K t13-08-1 p
8. --- Ho chiesto a un giovane \VUgras{garzone} \textsuperscript{(54)} che portava
\VUgras{bicchieri di birra} \textsuperscript{(55)} se la tavola rotonda fosse sotto la
veranda, e siccome mi ha detto di sì, mi sono seduto a uno dei tavoli e ho mangiato
benissimo.

%%K t13-noto-2 noto
\VUnotes{143.2mm}{%
\textit{\VUgras{(*) Con l'aiuto dell'elenco delle vivande dato nel vocabolario, il
maestro può variare facilmente il menu qui sotto.}}%
}

%%K t13-tit-2 sub
\VUsaut{4.0mm}
\VUpk{10.2pt}{40}{Un ottimo pranzo.}
\VUsaut{3.0mm}

%%K t13-09-1 p
9. --- Il cameriere mi ha portato la lista del giorno (*) e la lista dei vini. Ho
scelto e ho ordinato \VUgras{gamberetti} con \VUgras{burro} come
\VUgras{antipasto} e, come \VUgras{primo}, \VUgras{trippa alla moda di Caen}. Di
\VUgras{pesce} non ne ho preso, ma ho chiesto una porzione di \VUgras{cosciotto}
d'agnello e, come \VUgras{contorno}, \VUgras{spinaci alla panna}. Poi, siccome
nessuno dei \VUgras{dolci} mi attirava, mi sono fatto portare un po' di tutto:
\VUgras{formaggio}, \VUgras{frutta} e \VUgras{biscotti}. Ci ho aggiunto un ottimo
\VUgras{vino} Saint-Émilion e, molto contento del mio pranzo, mi sono alzato da tavola
dopo aver lasciato al \VUgras{cameriere} \textsuperscript{(57)} una bella
\VUgras{mancia} \textsuperscript{(59)}.

%%K t13-10-1 p
10. --- Vedendo che stavo per andarmene, il \VUgras{direttore} \textsuperscript{(60)}
mi ha proposto di uscire sul balcone ad ammirare lo splendido panorama
dell'\VUgras{oceano} \textsuperscript{(62)}. In lontananza si distinguevano piccole
\VUgras{barche} da pesca \textsuperscript{(63)}, \VUgras{velieri}
\textsuperscript{(64)} e un enorme \VUgras{transatlantico} \textsuperscript{(65)} la
cui sagoma nera si stagliava sulle \VUgras{nuvole} bianche \textsuperscript{(67)}
all'\VUgras{orizzonte} \textsuperscript{(66)}. Una brezza leggera muoveva la
\VUgras{bandiera} \textsuperscript{(70)} piantata sopra l'\VUgras{insegna}
\textsuperscript{(69)} del \VUgras{salone} \textsuperscript{(68)}.

%%K t13-tit-3 sub
\VUsaut{3.0mm}
\VUpk{10.2pt}{40}{Divertimenti.}
\VUsaut{3.0mm}

%%K t13-11-1 p
11. --- Lo spettacolo era così bello che ho deciso di salire il giorno dopo sulla
terrazza del caffè, portandomi un \VUgras{binocolo da teatro} \textsuperscript{(71)} o
un \VUgras{cannocchiale} \textsuperscript{(72)}.

%%K t13-12-1 p
12. --- Uscendo dal balcone sono andato al caffè dell'albergo per una
\VUgras{bibita} \textsuperscript{(73)} e una partita a \VUgras{biliardo}
\textsuperscript{(75)}. Ho attraversato la sala del caffè, dove molti clienti giocavano
a \VUgras{scacchi} \textsuperscript{(78)}, a \VUgras{dama} \textsuperscript{(79)}, a
\VUgras{domino} \textsuperscript{(80)}, a \VUgras{tavola reale} \textsuperscript{(81)} o
a \VUgras{carte} \textsuperscript{(82)}, e sono salito dritto alla \VUgras{sala del
biliardo} \textsuperscript{(74)}.

%%K t13-13-1 p
13. --- Finita la partita, sono sceso al pianterreno e ho chiesto al cameriere una
\VUgras{busta} \textsuperscript{(84)}, un po' di \VUgras{carta}
\textsuperscript{(83)}, un \VUgras{francobollo} \textsuperscript{(85)}, una
\VUgras{penna} \textsuperscript{(87)} e un po' d'\VUgras{inchiostro}
\textsuperscript{(86)} per scrivere una lettera.

%%K t13-13-2 p
Poi ho chiamato un \VUgras{vetturino} \textsuperscript{(92)} la cui \VUgras{carrozza}
\textsuperscript{(88)} si era fermata davanti al caffè. Gli ho chiesto il prezzo a ore
e a corsa, e quando ci siamo messi d'accordo mi ha aperto lo \VUgras{sportello}
\textsuperscript{(89)} e mi ha fatto salire. È risalito sulla sua \VUgras{cassetta}
\textsuperscript{(91)}, ha messo in moto i cavalli con la sua \VUgras{frusta}
\textsuperscript{(93)}, e via per una bellissima passeggiata.
\VUsaut{6.0mm}
\VUfilet{22mm}
